\section*{Introduction}

The same stimulus delivered to cortex rarely evokes the same response twice. Evoked activity adds to, and is scaled by,
ongoing fluctuations present at the moment of stimulation~\citep{Arieli1996, Petersen2003, Goris2014}. In awake mice
these fluctuations track movement and arousal~\citep{Niell2010, McGinley2015, Vinck2015, musall2019single,
stringer2019spontaneous, Salkoff2020}, alternate with slow, synchronized rhythms during quiet
wakefulness~\citep{Crochet2006, Poulet2008, Reimer2014}, and are shared across the two
hemispheres~\citep{Mohajerani2010, Shimaoka2019, Ye2026}. Responses to optogenetic and magnetic stimulation also change
with brain state~\citep{Mateo2011, Massimini2005, bryan2025temporal}, so a fixed stimulation protocol delivers a
different effective perturbation on every trial, which limits both the precision of causal manipulations and their
interpretation~\citep{Harris2011, Wolff2018}.

At the mesoscale, these fluctuations are directly observable. Wide-field calcium imaging resolves population activity
across the dorsal cortex of behaving mice~\citep{allen2017global, makino2017transformation, musall2019single,
Peters2021}, and has been combined with spatially targeted optogenetic stimulation to map and silence cortical
areas~\citep{Lim2012, ZatkaHaas2021}, now in awake mice over months~\citep{Matveev2024}. These experiments have been
open loop: the stimulus is fixed in advance, and the variability it meets is averaged away afterwards.

One way to tame this variability is to model it. Switching linear dynamical systems treat brain state as a parameter
that changes the dynamics~\citep{linderman2017rslds}, other approaches isolate the dynamics relevant to a
behavior~\citep{sani2021modeling}, and cortical state can now be decoded in real time~\citep{weiss2024realtime}. Yet
cortical activity also follows uninstructed movements and drive from other areas that a model of one region does not
observe~\citep{musall2019single, stringer2019spontaneous, zagha2013motor}, and univariate forecasts of spontaneous
wide-field activity stay informative for only about 1.5\,s~\citep{lu2025benchmarking}. A stimulation protocol that
anticipated each trial's state would need a model that identifies and forecasts it precisely.

Feedback needs much less. A controller that adjusts its input from a real-time measurement compensates for disturbances
it does not model~\citep{astrom1995pid}, so state-dependent variability becomes a disturbance to reject rather than a
confound to average~\citep{grosenick2015closed, RittChing2015}. Integral action, an internal model of a constant
signal, drives a constant offset toward zero~\citep{francis1976internal}, and at low gain it does so for an
exponentially stable linear plant despite actuator saturation and input delay, given only the sign of the plant's
steady-state gain~\citep{logemann2000absolute}. Closed-loop optogenetic control has held the firing rates of neurons at a
target in culture and in rodent thalamus~\citep{newman2015optogenetic, bolus2018design, bolus2021state}, clamped the
calcium signals of single neurons~\citep{Zhang2018}, and timed light to ongoing oscillations in rodents and
primates~\citep{Siegle2014, nicholson2018analogue, kanta2019closed, Zaaimi2022}. Input--output models of multiregional
responses to electrical stimulation have been proposed for model-based control, so far tested in closed loop only in
simulation~\citep{Yang2018, Yang2021}. At the mesoscale, real-time wide-field signals have closed the loop only through
behavior, as neurofeedback that animals learn to modulate~\citep{clancy2014volitional, Neely2018, Gupta2026}.

Two questions decide whether closed-loop control of mesoscale cortical activity is feasible and what it can reveal.
First, is the response to stimulation simple enough to design a controller around? Population dynamics are often well
captured by low-dimensional, linear descriptions~\citep{Churchland2012, vyas2020computation}, and linear models predict
resting-state human brain activity at the macroscale as well as nonlinear ones~\citep{Nozari2024}, but whether
optogenetic drive of a mesoscale signal is monotonic and approximately linear on average has not been tested. Second,
what does the error left by feedback reveal? A controller designed on the average response is not designed for
single-trial deviations, so relating the error it leaves to measured state can indicate which fluctuations lie beyond
its reach.

Here we establish closed-loop optogenetic control of mesoscale cortical activity in awake, head-fixed mice. We read out
regional $\Delta F/F$ in real time with wide-field calcium imaging and used it to set the power of inhibitory
optogenetic stimulation through a proportional--integral controller with feedforward compensation. Trial-averaged
suppression grew monotonically with stimulation power and was predicted by a low-order linear model, whereas single
trials varied with brain state about as much as ongoing activity did without stimulation. Across 15 sessions in four
mice, closed-loop control tracked a target activity level more accurately than open-loop stimulation in every session,
and with lower trial-to-trial variability on average. Against a disturbance predicted from the contralateral
hemisphere, feedback lowered residual error in 11 sessions, mostly by removing the offset left by open-loop
calibration. The error it left was larger on trials with more 1--4\,Hz power around stimulation, but not on trials with
more movement. In one mouse, adding a preview of the reference cancelled the phase lag in tracking a sinusoidal target,
and in a simulated replay of one session a model-predictive controller lowered the error left by feedback. These results
outline a procedure that does not depend on the readout or the stimulus: characterize the input--output relationship,
check that the average response is monotonic and approximately linear, and tune a proportional--integral controller to
the experiment. The error such a controller leaves then indicates which components of ongoing variability lie beyond a
fixed linear design, a starting point for context-dependent models and controllers.
